\documentclass[full-course-notes.tex]{subfiles}

\begin{document}
\lecturemeta{3}
  {Lexical Analysis: Tokens, Patterns, and Lexemes; the Chomsky Hierarchy; Regular Languages}
  {CLO 1}
  {\S3.1 (The Role of the Lexical Analyzer), \S3.3 (Specification of Tokens: terminology,
   strings, languages, operations)}
  {[AP] Modern Compiler Implementation, Ch.\,2}
\lecshort{Lexical Analysis \& Regular Languages}
\setcounter{section}{0}
\makelecturetitle

\begin{coreidea}[Lecture Roadmap]
\begin{enumerate}
  \item Where lexical analysis sits, and exactly how it talks to the parser; why lexer and
  parser are separate phases.
  \item The three core vocabulary words of this unit: \textbf{token}, \textbf{pattern},
  \textbf{lexeme} -- with several fully worked examples, plus token \textbf{attributes} and
  \textbf{lexical errors}, including a close look at \textbf{panic-mode recovery}.
  \item \textbf{The big picture}: how a regular expression eventually becomes a working,
  table-driven lexer, via finite automata -- the roadmap for this lecture and the ones after it.
  \item The \textbf{Chomsky hierarchy}: how formal, regular, context-free, and more powerful
  language classes relate to each other, and which machines recognize each one.
  \item \textbf{Regular languages}, defined formally and precisely -- the last topic of today,
  and the foundation the next lecture builds \textbf{regular expressions} on top of.
\end{enumerate}
By the end of this lecture you will know exactly which languages qualify as ``regular'' and
why, and you will have the full roadmap for how such a language eventually becomes a running
lexer. Lecture~\ref{lec:4} picks up immediately from here: regular expressions, the notation we
use to \emph{write down} a regular language.
\end{coreidea}

\section{Where Lexical Analysis Fits}
\label{sec:l3-where-lexical-fits}

\dbtag The \textbf{lexical analyzer} (also called a \textbf{scanner} or \textbf{lexer}) is the
first phase of a compiler. It reads the raw stream of characters making up the source program
and groups them into a stream of \textbf{tokens} --- meaningful chunks like keywords,
identifiers, numbers, and operators --- that the rest of the compiler works with instead of raw
text.

\begin{center}
\begin{tikzpicture}[
  box/.style={draw, thick, rounded corners, minimum height=1.1cm, minimum width=3cm, align=center, font=\small},
  tool/.style={box, fill=dbGreenBg, draw=dbGreen},
  data/.style={box, fill=nuLight, draw=nuNavy},
  >=Latex
]
  \node[data] (src) at (0,0) {Source Program\\(character stream)};
  \node[tool] (lex) at (4.2,0) {Lexical\\Analyzer};
  \node[tool] (parser) at (8.6,0) {Parser};
  \node[data] (st) at (4.2,-2.2) {Symbol Table};

  \draw[->, thick] (src) -- (lex);
  \draw[->, thick] ([yshift=2mm]lex.east) -- node[above, font=\scriptsize]{token} ([yshift=2mm]parser.west);
  \draw[->, thick] ([yshift=-2mm]parser.west) -- node[below, font=\scriptsize]{getNextToken()} ([yshift=-2mm]lex.east);
  \draw[<->, thick, dashed] (lex) -- (st);
\end{tikzpicture}
\end{center}

\begin{coreidea}[How the Lexer and Parser Actually Talk to Each Other]
In nearly every real compiler, the lexer does \emph{not} tokenize the whole file up front and
hand the parser a list. Instead, the parser drives the process: whenever the parser needs another
token, it calls a function conventionally named \texttt{getNextToken()}. The lexer resumes
scanning exactly where it left off, produces the next single token, and returns it. This
``pull'' model, sometimes called treating the lexer as a \textbf{coroutine} of the parser, avoids
ever materializing the full token list in memory and lets the lexer and parser interleave their
work one token at a time.
\end{coreidea}

\subsection{Why Not Just One Phase?}
\label{sec:l3-why-separate}

\dbtag DB gives three concrete engineering reasons for splitting lexical analysis out as its own
phase rather than folding it into the parser:

\begin{enumerate}
  \item \textbf{Simplicity of design.} Separating out low-level tasks (recognizing keywords,
  numbers, string literals) lets the parser work purely with grammatical structure, using
  tokens as its atomic symbols, instead of also worrying about individual characters.
  \item \textbf{Improved efficiency.} Specialized, table-driven techniques for recognizing
  tokens -- built from finite automata, which the lectures right after this one construct
  \emph{directly} from the regular expressions we develop later in this very lecture -- are
  considerably faster than the general-purpose parsing techniques would be if forced to operate
  character-by-character.
  \item \textbf{Enhanced portability.} Input-device and operating-system-specific quirks ---
  end-of-line conventions (\texttt{\textbackslash n} vs.\ \texttt{\textbackslash r\textbackslash n}), character
  encodings, special end-of-file characters --- are confined entirely to the lexer. The parser,
  and everything after it, is completely insulated from these differences.
\end{enumerate}

\begin{examfocus}
This ``three reasons'' list (simplicity, efficiency, portability) is a favorite short-answer
question. Be able to name and briefly justify all three, not just one.
\end{examfocus}

\section{Tokens, Patterns, and Lexemes}
\label{sec:l3-tokens-patterns-lexemes}

\dbtag These three terms are used precisely and are easy to blur together. Getting the
distinction exactly right is one of the most exam-tested ideas in this entire unit.

\begin{defbox}{Token, Pattern, Lexeme}
\begin{itemize}
  \item A \textbf{token} is a pair \texttt{<token-name, attribute-value>} representing a
  \emph{category} of lexical unit, e.g.\ \texttt{ID}, \texttt{NUM}, \texttt{RELOP}. The
  token-name is an abstract symbol used by the parser.
  \item A \textbf{pattern} is the rule that describes the \emph{set of possible lexemes} that
  can produce a given token -- for identifiers, informally, ``a letter followed by any number of
  letters and digits.'' Lecture~\ref{lec:4} makes patterns precise using regular expressions.
  \item A \textbf{lexeme} is the actual sequence of characters in the source program that
  \emph{matches} a pattern and is classified as an instance of some token.
\end{itemize}
\end{defbox}

\begin{center}
\begin{tikzpicture}[
  box/.style={draw, thick, rounded corners, minimum height=1cm, minimum width=3.1cm, align=center, font=\small},
  >=Latex
]
  \node[box, fill=nuLight, draw=nuNavy] (chars) at (0,0) {Source\\Characters};
  \node[box, fill=enrichGoldBg, draw=enrichGold] (pat) at (4,0) {matched against\\a \textbf{Pattern}};
  \node[box, fill=dbGreenBg, draw=dbGreen] (lex) at (8,0) {yields a\\\textbf{Lexeme}};
  \node[box, fill=examRedBg, draw=examRed] (tok) at (12,0) {emits a\\\textbf{Token}\\\texttt{<name, attr>}};
  \draw[->, thick] (chars) -- (pat);
  \draw[->, thick] (pat) -- (lex);
  \draw[->, thick] (lex) -- (tok);
\end{tikzpicture}
\end{center}

\begin{examplebox}{DB-Style Token/Pattern/Lexeme Table}
\renewcommand{\arraystretch}{1.3}
\begin{tabularx}{\textwidth}{@{}l X l@{}}
\toprule
\textbf{Token} & \textbf{Informal Description of Pattern} & \textbf{Sample Lexemes} \\
\midrule
\texttt{WHILE} & the exact keyword \texttt{while} & \texttt{while} \\
\texttt{ID} & a letter, followed by any number of letters or digits & \texttt{i}, \texttt{sum}, \texttt{score2} \\
\texttt{NUM} & any numeric constant & \texttt{0}, \texttt{42}, \texttt{3.14} \\
\texttt{RELOP} & \texttt{<}, \texttt{<=}, \texttt{=}, \texttt{<>}, \texttt{>}, or \texttt{>=} & \texttt{<=}, \texttt{<>} \\
\texttt{LITERAL} & any characters between a pair of double quotes, except a double quote itself & \texttt{"Total = "} \\
\bottomrule
\end{tabularx}
\end{examplebox}

\subsection{Worked Example 1: Tokenizing a Function Call}
\label{sec:l3-worked-example-1}

Consider the following line of C-like source code:

\begin{center}
\ttfamily\large
\colorbox{dbGreenBg}{printf}\colorbox{examRedBg}{(}\colorbox{enrichGoldBg}{"Total = \%d\textbackslash n"}\colorbox{examRedBg}{,}\ \colorbox{dbGreenBg}{score}\colorbox{examRedBg}{)}\colorbox{examRedBg}{;}
\end{center}
\begin{center}
\small\begin{tabular}{@{}p{2.6cm}p{2.6cm}p{2.6cm}p{2.6cm}@{}}
\colorbox{dbGreenBg}{\phantom{x}}\ identifier &
\colorbox{examRedBg}{\phantom{x}}\ punctuation/operator &
\colorbox{enrichGoldBg}{\phantom{x}}\ string literal &
\end{tabular}
\end{center}

\begin{center}
\renewcommand{\arraystretch}{1.25}
\begin{tabularx}{\textwidth}{@{}l l X@{}}
\toprule
\textbf{Lexeme} & \textbf{Token Name} & \textbf{Attribute-Value (passed to parser)} \\
\midrule
\texttt{printf} & \texttt{ID} & pointer to symbol-table entry for \texttt{printf} \\
\texttt{(} & \texttt{LPAREN} & none needed \\
\texttt{"Total = \%d\textbackslash n"} & \texttt{LITERAL} & pointer to the string's stored value \\
\texttt{,} & \texttt{COMMA} & none needed \\
\texttt{score} & \texttt{ID} & pointer to symbol-table entry for \texttt{score} \\
\texttt{)} & \texttt{RPAREN} & none needed \\
\texttt{;} & \texttt{SEMI} & none needed \\
\bottomrule
\end{tabularx}
\end{center}

\section{Attributes for Tokens}
\label{sec:l3-attributes}

\dbtag When more than one lexeme can match a token's pattern (as with \texttt{ID} or
\texttt{NUM}), the lexer must pass along enough extra information -- the \textbf{attribute-value}
-- for later phases to know exactly \emph{which} identifier or \emph{which} number was matched.
For identifiers, this is essentially always a pointer into the symbol table.

\begin{examplebox}{Worked Example 2: An Assignment Statement}
Source line: \quad \texttt{E = M * C ** 2}

Assume \texttt{**} is this language's exponentiation operator. The lexer's output, as a stream
of \texttt{<token-name, attribute-value>} pairs, is:
\begin{center}
\ttfamily
\begin{tabular}{@{}l@{}}
\ttfamily <ID, ptr to symtab entry for E> \\
<ASSIGN\_OP,\ -- >  \\
<ID, ptr to symtab entry for M> \\
<MULT\_OP,\ -- > \\
<ID, ptr to symtab entry for C> \\
<EXP\_OP,\ -- > \\
<NUM, integer value 2>
\end{tabular}
\end{center}
Notice: \texttt{ID} tokens all carry a symbol-table pointer as their attribute so that later
phases can distinguish \texttt{E} from \texttt{M} from \texttt{C}; the operator tokens need no
attribute at all, since the token-name alone (\texttt{ASSIGN\_OP} vs.\ \texttt{MULT\_OP}) already
says everything the parser needs to know.
\end{examplebox}

\begin{examfocus}
Given a short program line, be able to produce the exact stream of
\texttt{<token-name, attribute-value>} pairs, in order, the way both worked examples above do.
This is the single most common long-answer question style for this unit.
\end{examfocus}

\subsection{Worked Example 3: A Complete Statement Block}
\label{sec:l3-worked-example-3}

\begin{center}
\begin{minipage}{0.5\textwidth}
\ttfamily
int sum = 0;\\
while (i <= 10) \{\\
\ \ sum = sum + i;\\
\ \ i = i + 1;\\
\}
\end{minipage}
\end{center}

\begin{center}
\renewcommand{\arraystretch}{1.15}
\footnotesize
\begin{tabularx}{\textwidth}{@{}l l X | l l X@{}}
\toprule
\textbf{Lexeme} & \textbf{Token} & \textbf{Attribute} & \textbf{Lexeme} & \textbf{Token} & \textbf{Attribute} \\
\midrule
\texttt{int} & \texttt{TYPE} & -- & \texttt{sum} & \texttt{ID} & ptr(\texttt{sum}) \\
\texttt{sum} & \texttt{ID} & ptr(\texttt{sum}) & \texttt{=} & \texttt{ASSIGN} & -- \\
\texttt{=} & \texttt{ASSIGN} & -- & \texttt{sum} & \texttt{ID} & ptr(\texttt{sum}) \\
\texttt{0} & \texttt{NUM} & 0 & \texttt{+} & \texttt{PLUS} & -- \\
\texttt{;} & \texttt{SEMI} & -- & \texttt{i} & \texttt{ID} & ptr(\texttt{i}) \\
\texttt{while} & \texttt{WHILE} & -- & \texttt{;} & \texttt{SEMI} & -- \\
\texttt{(} & \texttt{LPAREN} & -- & \texttt{i} & \texttt{ID} & ptr(\texttt{i}) \\
\texttt{i} & \texttt{ID} & ptr(\texttt{i}) & \texttt{=} & \texttt{ASSIGN} & -- \\
\texttt{<=} & \texttt{RELOP} & LE & \texttt{i} & \texttt{ID} & ptr(\texttt{i}) \\
\texttt{10} & \texttt{NUM} & 10 & \texttt{+} & \texttt{PLUS} & -- \\
\texttt{)} & \texttt{RPAREN} & -- & \texttt{1} & \texttt{NUM} & 1 \\
\texttt{\{} & \texttt{LBRACE} & -- & \texttt{;} & \texttt{SEMI} & -- \\
 & & & \texttt{\}} & \texttt{RBRACE} & -- \\
\bottomrule
\end{tabularx}
\end{center}

Two details worth noticing in this longer example:
\begin{itemize}
  \item \texttt{sum} and \texttt{i} each get \emph{their own} symbol-table pointer the first
  time they're seen, and the \emph{same} pointer every time they reappear -- the symbol table
  guarantees one entry per distinct identifier, no matter how many times it's used.
  \item Whitespace (the spaces and newlines between tokens) and indentation produce \emph{no}
  tokens at all -- they are consumed and discarded by the lexer, along with any comments, before
  the parser ever sees anything.
\end{itemize}

\begin{coreidea}[The Lexer's ``Housekeeping'' Duties]
Beyond producing tokens, DB notes the lexer is also the natural place to: strip out whitespace
and comments (neither produces a token); and track line numbers (and often column numbers) so
that later phases can report errors with accurate source locations, e.g.\ ``\texttt{undeclared
variable `x' at line 42}.''
\end{coreidea}

\section{Lexical Errors}
\label{sec:l3-lexical-errors}

\dbtag It is surprisingly rare for a lexer to be able to detect an error entirely on its own,
\emph{by itself, in isolation} -- because almost any short character sequence is a \emph{prefix}
of some valid lexeme for some token.

\begin{examplebox}{The Classic Illustration}
Consider the fragment \texttt{fi (a == f(x))} in a C-like language, where the programmer
\emph{meant} to type the keyword \texttt{if}. The lexer cannot know this. \texttt{fi} is a
perfectly legal lexeme for token \texttt{ID} (``a letter followed by letters/digits''), so the
lexer happily emits \texttt{<ID, ptr(fi)>}. The mistake will only be caught later -- by the
\emph{parser}, if \texttt{fi} appears somewhere no identifier is grammatically valid, or by
\emph{semantic analysis}, if \texttt{fi} is used as if it were a function without ever having
been declared.
\end{examplebox}

Genuine lexical errors are ones where \emph{no pattern for any token} can match the remaining
input at all, or where a token's own internal structure is malformed:

\begin{itemize}
  \item An illegal character that starts no valid token in the language, e.g.\ a stray
  \texttt{\#} or \texttt{@} in a language that assigns those characters no meaning.
  \item An unterminated string literal: opening quote seen, but the line (or file) ends before a
  matching closing quote.
  \item A malformed numeric literal, e.g.\ \texttt{3.14.15} (two decimal points), or
  \texttt{2r} immediately followed by a letter with no valid meaning in that position.
  \item An unterminated comment: a \texttt{/*} seen with no matching \texttt{*/} anywhere before
  end of file.
\end{itemize}

\subsection{Error Recovery Strategies}
\label{sec:l3-error-recovery}

\dbtag A production-quality lexer does not simply stop at the first error --- it tries to
recover and keep going, so the programmer can see \emph{several} distinct errors from one
compilation attempt rather than fixing them one at a time.

\begin{center}
\begin{tikzpicture}[
  box/.style={draw, thick, rounded corners, minimum height=0.9cm, minimum width=3.6cm, align=center, font=\small},
  dec/.style={draw, thick, diamond, aspect=2.4, minimum width=3.4cm, minimum height=1.4cm, align=center, font=\small, fill=nuLight, draw=nuNavy},
  >=Latex
]
  \node[box, fill=nuLight, draw=nuNavy] (start) at (0,0) {Read next\\characters};
  \node[dec] (match) at (0,-2.1) {Matches some\\token's pattern?};
  \node[box, fill=dbGreenBg, draw=dbGreen] (emit) at (5,-2.1) {Emit token,\\advance input};
  \node[box, fill=examRedBg, draw=examRed] (err) at (0,-4.6) {Report lexical\\error};
  \node[box, fill=enrichGoldBg, draw=enrichGold] (panic) at (0,-6.6) {Panic-mode: delete\\char(s) until a valid\\token can resume};

  \draw[->, thick] (start) -- (match);
  \draw[->, thick] (match) -- node[above, font=\scriptsize]{yes} (emit);
  \draw[->, thick] (match) -- node[left, font=\scriptsize]{no} (err);
  \draw[->, thick] (err) -- (panic);
  \draw[->, thick] (emit.north) to[out=90,in=0] (start.east);
  \draw[->, thick] (panic.west) to[out=180,in=270] (start.south);
\end{tikzpicture}
\end{center}

\dbtag \textbf{Panic mode, in detail.} ``Panic mode'' names a very specific, very simple
procedure, and it helps to spell out its steps exactly, since the name alone makes it sound more
dramatic than it is: (1)~report the error at the current position; (2)~delete (discard) the
current character; (3)~check whether the \emph{remaining} input now begins with something
matching a valid token's pattern; (4)~if not, go back to step 2 and delete the next character
too; (5)~as soon as the remaining input does match a valid token -- or a natural boundary like
end of line or end of file is reached -- stop deleting and resume ordinary scanning from there.
In short: \emph{throw away characters, one at a time, until scanning can safely restart.}

\begin{examplebox}{Worked Example 1: A Single Illegal Character}
Source line: \quad \texttt{total = x @ y;}

The lexer scans \texttt{total} (\texttt{ID}), \texttt{=} (\texttt{ASSIGN}), \texttt{x}
(\texttt{ID}) normally. It then reaches \texttt{@}, which matches no token pattern in this
language. \textbf{Step 1:} report ``illegal character \texttt{@} at column 12.'' \textbf{Steps
2--3:} delete just the \texttt{@}; the very next character, a space, already leads into a
lexeme (\texttt{y}) that matches \texttt{ID}. So panic mode stops after deleting a single
character. Scanning resumes normally: \texttt{y} (\texttt{ID}), \texttt{;} (\texttt{SEMI}). One
error reported, and the rest of the line recovers cleanly -- this is the best case for panic
mode.
\end{examplebox}

\begin{examplebox}{Worked Example 2: Discarding Until a Natural Boundary}
Source:
\begin{verbatim}
printf("Hello world);
x = 5;
\end{verbatim}
The lexer scans \texttt{printf} (\texttt{ID}), \texttt{(} (\texttt{LPAREN}), then starts
matching a \texttt{LITERAL} token at the opening \texttt{"}. It looks for a matching closing
\texttt{"} but reaches the end of the line without finding one. Report: ``unterminated string
literal starting at line 1.'' Unlike Worked Example 1, there is no small number of characters
to delete that fixes this -- \emph{nothing} inside a broken string is a safe resumption point.
Panic mode instead discards everything from the opening quote through the end of the line (a
natural boundary most languages use for exactly this reason: string literals cannot normally
span a raw newline), and resumes scanning fresh at line 2: \texttt{x} (\texttt{ID}), \texttt{=}
(\texttt{ASSIGN}), \texttt{5} (\texttt{NUM}), \texttt{;} (\texttt{SEMI}). One error reported;
the rest of the program is still checked for further problems.
\end{examplebox}

\begin{examplebox}{Worked Example 3: Why Panic Mode Can Cascade Into Extra Errors}
Source line: \quad \texttt{y = 3.14.159 + 2;}

The lexer scans \texttt{y} (\texttt{ID}), \texttt{=} (\texttt{ASSIGN}), then starts matching a
\texttt{NUM} at \texttt{3.14}. Maximal munch tries to extend the match, but a \emph{second}
\texttt{.} cannot extend a valid \texttt{num} (DB's regular definition, covered fully with
Lecture~\ref{lec:4}'s \texttt{num} example, allows only one optional fraction part) -- so the lexer accepts what it \emph{did}
successfully match, emits \texttt{<NUM, 3.14>}, and resumes scanning from the very next
character: a bare \texttt{.}. But a lone \texttt{.} does not start any valid token in this
language either, so panic mode fires a \emph{second} time, reporting a second illegal-character
error and deleting just that \texttt{.}, before finally matching \texttt{159} as another
\texttt{NUM} token. \textbf{One root mistake -- a stray decimal point -- surfaces as two
reported errors.} This is a known, unavoidable downside of panic mode: it recovers
\emph{syntactically} (scanning can always continue), not \emph{semantically} (it has no idea
what the programmer actually meant), so one typo can occasionally produce a short burst of
related-looking error messages. Experienced programmers learn to fix the \emph{first} error in
a cluster and recompile, rather than trying to fix every reported line individually.
\end{examplebox}

\dbtag \textbf{Beyond panic mode: local corrections.} DB also lists four more surgical
``local-correction'' strategies a recovering lexer \emph{could} attempt instead of blindly
discarding characters -- though, as the note after the examples explains, only two of them
apply cleanly to genuine lexical errors:

\begin{itemize}
  \item \textbf{Deleting} an extraneous character -- e.g.\ the malformed number
  \texttt{3..14} (one \texttt{.} too many): deleting the second \texttt{.} recovers the valid
  number \texttt{3.14}, instead of panic mode's blunter ``discard until something matches.''
  \item \textbf{Inserting} a missing character -- e.g.\ the unterminated string \texttt{"Hello}
  with the line ending before a closing quote: inserting a \texttt{"} right where the line ends
  recovers a (short) valid string literal, letting scanning continue past it immediately
  instead of discarding the rest of the line as Worked Example 2 did.
  \item \textbf{Replacing} one character with another -- e.g.\ \texttt{3@4} plausibly intended
  as \texttt{3*4}: replacing \texttt{@} with \texttt{*} recovers a valid expression.
  \item \textbf{Transposing} two adjacent characters -- e.g.\ \texttt{"Helol"} plausibly
  intended as \texttt{"Hello"}.
\end{itemize}

\begin{examfocus}
Not all four local corrections are equally ``lexical.'' \textbf{Deletion} and \textbf{insertion}
directly repair the two genuine categories of lexical error from
\Sref{sec:l3-lexical-errors} -- an extra character the pattern can't accommodate, or a
character the pattern is missing -- so a lexer's own panic-mode-plus-local-correction logic can
apply them without knowing anything about programmer intent. \textbf{Replacement} and
\textbf{transposition}, by contrast, mostly turn one \emph{valid} lexeme into a
\emph{different} valid lexeme (as the \texttt{Helol}$\to$\texttt{Hello} example shows) --
exactly the ``looks like a typo but is a perfectly legal token'' situation from the
\texttt{fi}/\texttt{if} example earlier in this lecture, which is \emph{not} a lexical error at
all and is therefore outside what the lexer itself can or should silently fix. Replacement and
transposition matter far more for IDE-style ``did you mean \dots?'' suggestions layered on top
of the compiler than for the lexer's own error-recovery logic.
\end{examfocus}

\begin{historynote}
DB points toward a more principled version of this idea: choosing whichever single correction
requires the \emph{fewest} character edits -- insertions, deletions, substitutions -- to turn
the bad input into something valid. This is exactly \textbf{edit distance} (Levenshtein
distance), the same string-similarity measure used today in spell-checkers and DNA sequence
alignment. Aho and Peterson described a minimum-distance error-correcting parser as early as
1972; the idea that ``the best fix is the cheapest fix'' turns out to generalize far beyond
lexical analysis, into syntax error recovery and even runtime auto-correct systems.
\end{historynote}

\begin{examfocus}
Be able to explain, with a concrete example like \texttt{fi (a == f(x))}, \emph{why} a
misspelled keyword is usually \textbf{not} a lexical error even though it looks like an obvious
typo. Then be able to walk through panic mode step by step (report, delete, check, repeat,
resume) on a short broken input the way the three worked examples above do -- this is the most
common long-answer version of this topic, and students often lose marks by describing panic
mode too vaguely (``it skips bad characters'') instead of naming the actual step-by-step
procedure.
\end{examfocus}

\section{The Big Picture: From Regular Expressions to a Working Lexer}
\label{sec:l3-big-picture}

\dbtag Every pattern we have used so far has been informal English: ``a letter followed by
letters or digits.'' That is fine for talking about lexical analysis, but it is not something a
computer can execute. Turning informal patterns into an actual running lexer takes a short
pipeline of well-defined, \emph{mechanical} steps, and it is worth seeing the whole pipeline
once, up front, before building any one piece of it.

\begin{center}
\resizebox{\textwidth}{!}{%
\begin{tikzpicture}[
  box/.style={draw, thick, rounded corners, minimum height=1.1cm, minimum width=2.7cm, align=center, font=\small},
  re/.style={box, fill=examRedBg, draw=examRed},
  nfa/.style={box, fill=enrichGoldBg, draw=enrichGold},
  dfa/.style={box, fill=dbGreenBg, draw=dbGreen},
  tab/.style={box, fill=nuLight, draw=nuNavy},
  >=Latex
]
  \node[re]  (re)    at (0,0)    {Regular\\Expression};
  \node[nfa] (nfa)   at (4.6,0)  {NFA};
  \node[dfa] (dfa)   at (9.2,0)  {DFA};
  \node[dfa] (mdfa)  at (13.8,0) {Minimized\\DFA};
  \node[tab] (table) at (18.4,0) {Transition\\Table /\\Diagram};
  \node[tab] (lexer) at (18.4,-2.6) {Table-Driven\\Lexer};

  \draw[->, thick] (re) -- node[above, font=\scriptsize, align=center]{Thompson's\\Construction} (nfa);
  \draw[->, thick] (nfa) -- node[above, font=\scriptsize, align=center]{Subset\\Construction} (dfa);
  \draw[->, thick] (dfa) -- node[above, font=\scriptsize, align=center]{Minimization} (mdfa);
  \draw[->, thick] (mdfa) -- node[above, font=\scriptsize, align=center]{Table\\Encoding} (table);
  \draw[->, thick] (table) -- (lexer);
  \draw[->, thick, dashed, black!55] (re.south) to[out=-55, in=235]
    node[below=3pt, font=\scriptsize, align=center, text=black!55]
    {\textbf{Shortcut:} direct DFA construction\\(skips the NFA step entirely)} (dfa.south);
\end{tikzpicture}}
\end{center}

\begin{center}
\renewcommand{\arraystretch}{1.35}
\begin{tabularx}{\textwidth}{@{}l X@{}}
\toprule
\textbf{Step} & \textbf{What Happens} \\
\midrule
Regular expression $\to$ NFA & \textbf{Thompson's construction}: a purely mechanical algorithm
that builds a nondeterministic finite automaton (NFA) directly from the structure of a regular
expression, one small automaton per operator. \\
NFA $\to$ DFA & \textbf{Subset construction}: another mechanical algorithm that converts any NFA
into an equivalent deterministic finite automaton (DFA) -- one with no ambiguity about which
transition to take, at the cost of (possibly) more states. \\
\emph{(shortcut)} Regular expression $\to$ DFA directly & Some tools skip the NFA step
entirely: DB \S3.9 builds a DFA \emph{directly} from a regular expression's syntax tree, using
functions traditionally called \textit{nullable}, \textit{firstpos}, \textit{lastpos}, and
\textit{followpos}. It computes exactly the same DFA that Thompson's construction plus subset
construction would have produced (possibly needing a separate minimization pass afterward) --
it is a different \emph{route} to the same destination, not a different destination, and it is
what many real lexer generators use internally because it tends to build fewer intermediate
states. \\
DFA $\to$ minimized DFA & \textbf{DFA minimization}: merges states that are behaviorally
indistinguishable, producing the smallest possible DFA that still recognizes exactly the same
language. \\
Minimized DFA $\to$ transition table & The (finite!) set of states and transitions is encoded
as a 2-D table: \texttt{table[state][input\_symbol] = next\_state}. \\
Transition table $\to$ lexer & A tiny, fast loop -- look up the current state and next
character in the table, move to the next state, repeat -- \emph{is} the lexer's inner loop. No
more cleverness is needed once the table exists. \\
\bottomrule
\end{tabularx}
\end{center}

\begin{examfocus}
``Transition diagram'' and ``transition table'' are not two different stages of this pipeline --
they are two different \emph{representations of the same DFA (or NFA)}. A \textbf{transition
diagram} is the picture: circles for states, labeled arrows for transitions, exactly the
automaton drawings you will see in the next few lectures. A \textbf{transition table} is the
same information laid out as a 2-D array instead of a picture. Neither one is more ``correct''
than the other, and neither is a separate construction step -- once you have a DFA, you already
have both, you have just chosen how to write it down. A common exam trap asks you to convert
between the two for a small automaton; that is a notation exercise, not a new algorithm.
\end{examfocus}

\begin{coreidea}[Why This Pipeline Matters]
Every solid arrow in the diagram above is a \emph{mechanical, automatable} algorithm -- none of
them require human judgment once the previous step is done, and the dashed shortcut is simply
an alternative route to the same DFA. This is exactly why lexer-generator tools like Flex exist:
you write only the left-most box (a set of regular expressions, one per token), and the tool
runs the entire rest of the pipeline for you, emitting ready-to-compile scanner code. The
remaining lectures on lexical analysis build this pipeline one arrow at a time -- Thompson's
construction and subset construction first, then minimization -- so that by the time we reach
Flex itself, you understand exactly what the tool is doing on your behalf rather than treating
it as a black box.
\end{coreidea}

\section{Strings, Alphabets, and Languages}
\label{sec:l4-strings-alphabets-languages}

\dbtag Before regular expressions can be defined precisely, three primitive terms need to be
pinned down exactly as DB uses them.

\begin{defbox}{Alphabet, String, Language}
\begin{itemize}
  \item An \textbf{alphabet} $\Sigma$ is any finite set of symbols, e.g.\ $\{0,1\}$ or the set
  of ASCII characters.
  \item A \textbf{string} (or \emph{word}) over $\Sigma$ is a finite sequence of symbols drawn
  from $\Sigma$. The \textbf{empty string}, written $\varepsilon$, has length 0.
  \item A \textbf{language} over $\Sigma$ is simply \emph{any} set of strings over $\Sigma$ ---
  possibly infinite, possibly empty ($\emptyset$, the language with no strings at all, which is
  \emph{not} the same as $\{\varepsilon\}$, the language containing only the empty string).
\end{itemize}
\end{defbox}

\begin{examfocus}
$\emptyset$ and $\{\varepsilon\}$ are a classic confusion. $\emptyset$ contains \emph{zero}
strings; $\{\varepsilon\}$ contains \emph{one} string (the empty one). Keep them distinct.
\end{examfocus}

\begin{coreidea}[Formal Languages: A Hierarchy, Not One Thing]
The definition above -- ``a language is any set of strings'' -- deliberately puts \emph{no}
restriction on what that set looks like. A set meeting only this bare definition is often
called a \textbf{formal language}, precisely to flag that nothing about its internal structure
is assumed yet. Computer science studies several increasingly restrictive (and increasingly
cheap to process) \emph{classes} of formal languages, most famously organized into the
\textbf{Chomsky hierarchy} -- the subject of the subsection below.
\end{coreidea}

\subsection{The Chomsky Hierarchy}
\label{sec:l3-chomsky-hierarchy}

\dbtag Linguist Noam Chomsky proposed this hierarchy in 1956 while studying the structure of
natural language; computer science adopted it almost immediately, because it turns out to
classify exactly the right notion for programming languages too -- \emph{how much machinery do
you need to recognize this language?} Each class below is defined by a restriction on the
\emph{grammar} allowed to generate it, and each corresponds to exactly one class of recognizing
machine.

\begin{center}
\resizebox{0.62\textwidth}{!}{%
\begin{tikzpicture}
  \foreach \r/\fillcol/\drawcol in {%
    6.0/nuLight/nuNavy,%
    4.8/enrichGoldBg/enrichGold,%
    3.6/examRedBg/examRed,%
    2.4/white/gray,%
    1.2/dbGreenBg/dbGreen%
  } {
    \fill[\fillcol] (0,0) circle (\r);
    \draw[thick, \drawcol] (0,0) circle (\r);
  }
  \node[align=center, font=\small, text=nuNavy] at (0,5.1) {\textbf{Type 0: RE}\\(Turing Machine)};
  \node[align=center, font=\small, text=enrichGold] at (0,3.9) {\textbf{Type 1: CS}\\(LBA)};
  \node[align=center, font=\small, text=examRed] at (0,2.7) {\textbf{Type 2: CF}\\(NPDA)};
  \node[align=center, font=\small, text=gray] at (0,1.5) {\textbf{Det.\ CFL}\\(DPDA)};
  \node[align=center, font=\small, text=dbGreen] at (0,0) {\textbf{Type 3}\\(DFA=NFA)};
\end{tikzpicture}}
\end{center}
\begin{center}
\small\textit{RE = Recursively Enumerable, CS = Context-Sensitive, CF = Context-Free -- full
names and machines in the table below.}
\end{center}
\begin{center}
\small\textit{Each ring properly \emph{contains} the ones inside it: every regular language is
context-free, every context-free language is context-sensitive, and so on -- but not the
reverse.}
\end{center}

\begin{center}
\small
\renewcommand{\arraystretch}{1.5}
\begin{tabularx}{\textwidth}{@{}c l l X@{}}
\toprule
\textbf{Type} & \textbf{Class} & \textbf{Machine} & \textbf{Deterministic vs.\ Nondeterministic} \\
\midrule
3 & Regular & DFA or NFA & \textbf{Equal power.} Every NFA has an equivalent DFA (via the
\emph{subset construction} we build toward in the next few lectures) -- nondeterminism buys
convenience of expression here, not extra recognizing power. \\
-- & Deterministic CFL & DPDA & \textbf{Strictly weaker than NPDA.} Deterministic pushdown
automata recognize only a proper subset of the context-free languages -- this middle ring is
not one of Chomsky's original four types, but is a genuinely useful refinement between Types 2
and 3. \\
2 & Context-Free & PDA (nondeterministic, NPDA) & \textbf{Strictly stronger than DPDA.} The
classic separating example: $\{ww^{R} : w \in \{a,b\}^{*}\}$ (even-length palindromes) is
context-free, but no DPDA can recognize it -- a deterministic machine can't ``guess'' where the
middle of the string is without trying multiple possibilities at once. \\
1 & Context-Sensitive & Linear-Bounded Automaton (LBA) & \textbf{Unknown!} Whether
deterministic and nondeterministic LBAs have equal power is the \textbf{LBA problem}, open
since 1964 -- one of the oldest unsolved questions in theoretical computer science. \\
0 & Recursively Enumerable & Turing Machine & \textbf{Equal power, and this is the ceiling.}
Every reasonable model of computation ever proposed -- multi-tape TMs, nondeterministic TMs,
RAM machines, $\lambda$-calculus -- has been proven exactly as powerful as the single-tape
deterministic TM. This equivalence, and the belief that no more powerful model of
\emph{effective computation} exists, is the \textbf{Church--Turing thesis}. \\
\bottomrule
\end{tabularx}
\end{center}

\begin{examfocus}
Keep the \textbf{NFA/DFA} case and the \textbf{PDA/DPDA} case straight -- they go opposite
ways, and exams like to test exactly this contrast. For \emph{finite} automata, adding
nondeterminism costs nothing in power (only in state count). For \emph{pushdown} automata,
adding nondeterminism genuinely lets you recognize more languages. This single fact is the
formal reason a lexer can always be built deterministically (fast, table-driven), while a
general context-free parser cannot always be -- some grammars force real backtracking or lookahead
machinery that a DPDA-equivalent design cannot avoid.
\end{examfocus}

This lecture defines exactly one of these classes with full precision -- the \textbf{regular
languages}, starting below. \textbf{Context-free languages} are the parser's formalism, and the
course returns to them once we reach syntax analysis; every regular language is also
context-free (but not the reverse), which is exactly why a parser can express everything a
lexer's regular expressions can, and more.

\section{Regular Languages: A Formal Definition}
\label{sec:l3-regular-languages-formal}

\dbtag To say precisely which languages qualify as ``regular,'' we first need a small set of
operations for building new languages out of old ones. Given two languages $L$ and $M$ over the
same alphabet, DB defines four such operations:

\begin{center}
\renewcommand{\arraystretch}{1.35}
\begin{tabularx}{\textwidth}{@{}l l X@{}}
\toprule
\textbf{Operation} & \textbf{Notation} & \textbf{Definition} \\
\midrule
Union & $L \cup M$ & $\{\, s \mid s \in L \text{ or } s \in M \,\}$ \\
Concatenation & $LM$ & $\{\, st \mid s \in L \text{ and } t \in M \,\}$ (paste a string from $L$
directly onto a string from $M$) \\
Kleene closure & $L^{*}$ & $\bigcup_{i=0}^{\infty} L^{i}$ -- zero or more strings from $L$,
concatenated (includes $\varepsilon$, from $L^0$) \\
Positive closure & $L^{+}$ & $\bigcup_{i=1}^{\infty} L^{i}$ -- one or more strings from $L$,
concatenated (excludes $\varepsilon$ unless $\varepsilon \in L$) \\
\bottomrule
\end{tabularx}
\end{center}

\begin{examplebox}{Worked Example: Letters and Digits}
Let $L = \{a, b, c, \dots, z, A, \dots, Z\}$ (the 52 letters) and $D = \{0, 1, \dots, 9\}$ (the
10 digits), each treated as a language of one-character strings.
\begin{itemize}
  \item $L \cup D$ -- every single letter \emph{or} single digit: 62 strings, each of length 1.
  \item $LD$ -- every string formed by one letter followed by one digit, e.g.\ \texttt{a5},
  \texttt{Z0} -- exactly $52 \times 10 = 520$ strings, each of length 2.
  \item $L^{4}$ -- every string of exactly 4 letters, e.g.\ \texttt{Test}, \texttt{abcd} --
  there are $52^{4}$ such strings in total.
  \item $L^{*}$ -- every string of letters of \emph{any} length, including $\varepsilon$ ---
  infinitely many strings.
  \item $L(L \cup D)^{*}$ -- one letter, followed by zero or more letters-or-digits. This is
  \emph{exactly} the language of identifiers in most C-like languages! We will turn this
  directly into a regular expression in Lecture~\ref{lec:4}.
\end{itemize}
\end{examplebox}

\dbtag With these operations in hand, we can now say precisely what ``regular'' means -- a
statement about the language itself, as a set of strings, before we ever write a single regular
expression on paper.

\begin{defbox}{Regular Language (Inductive Definition)}
The \textbf{regular languages} over an alphabet $\Sigma$ are defined inductively, using exactly
the union, concatenation, and Kleene-closure operations above.

\textbf{Base cases.} Each of the following is a regular language:
\begin{enumerate}[label=(\alph*)]
  \item $\emptyset$ -- the empty language, containing no strings at all;
  \item $\{\varepsilon\}$ -- the language containing only the empty string;
  \item $\{a\}$, for every symbol $a \in \Sigma$ -- the language containing just that one
  single-character string.
\end{enumerate}

\textbf{Inductive cases.} If $L$ and $M$ are regular languages, then each of the following is
\emph{also} a regular language:
\begin{enumerate}[label=(\arabic*)]
  \item $L \cup M$ -- their union;
  \item $LM$ -- their concatenation;
  \item $L^{*}$ -- the Kleene closure of $L$.
\end{enumerate}

\textbf{Closure clause.} Nothing else is a regular language: the regular languages are exactly
the languages reachable by finitely many applications of the three inductive cases above,
starting only from the three base cases.
\end{defbox}

This definition tells us \emph{which} languages count as regular. It does \emph{not}, by
itself, give us a convenient way to write such a language down on paper or in code -- and that
gap is exactly what a regular expression fills.

\begin{coreidea}[Regular Language vs.\ Regular Expression]
A \textbf{regular language} is a set of strings satisfying the definition above -- a semantic
object, a mathematical fact about which strings belong to it. A \textbf{regular expression}
(next lecture) is a piece of \emph{syntax}: a compact, human-writable notation for naming one
specific regular language. Every regular expression denotes exactly one regular language, and
--- a real theorem, proved properly in your Automata Theory course --- every regular language
can be named by at least one regular expression. The two notions therefore describe exactly the
same class of languages, but they are conceptually different: one is a fact about a language,
the other is a way of writing it down.
\end{coreidea}

\begin{examfocus}
``Is this a regular language?'' and ``is this a valid regular expression?'' are different
questions, and exam wording sometimes blurs them on purpose. A language can be regular even if
you haven't yet written down a regular expression for it (you just haven't found the notation
yet); conversely, once we define regular expressions formally next lecture, every string of
regex syntax that type-checks will denote \emph{some} regular language, whether or not that
language is a useful one.
\end{examfocus}

\begin{takeaways}
\begin{itemize}
  \item A \textbf{lexeme} is matched source text; a \textbf{token} is the abstract
  \texttt{<name, attribute>} pair a lexer emits for it; a \textbf{pattern} is the rule
  connecting the two.
  \item Genuine \textbf{lexical errors} are things no pattern can match at all (illegal
  characters, unterminated strings/comments, malformed numbers) -- a misspelled keyword is
  \emph{not} one of these, since it's still a valid identifier lexeme.
  \item \textbf{Panic mode} reports an error, then discards characters one at a time until
  scanning can safely resume -- simple, always terminates, but can occasionally cascade into
  extra reported errors from one root mistake.
  \item The big picture: \textbf{regular expression $\to$ NFA $\to$ DFA $\to$ minimized DFA
  $\to$ transition table $\to$ lexer}, via Thompson's construction, subset construction, and
  minimization (or a direct regex-to-DFA shortcut) -- every step mechanical, which is exactly
  why lexer generators can automate all of it.
  \item The \textbf{Chomsky hierarchy} nests formal languages by how much machine power they
  need: regular (DFA = NFA) $\subset$ context-free (PDA, strictly more powerful
  nondeterministic than deterministic) $\subset$ context-sensitive (LBA) $\subset$ recursively
  enumerable (Turing machine, the most powerful model known).
  \item A \textbf{regular language} is a set of strings, defined inductively from three base
  cases ($\emptyset$, $\{\varepsilon\}$, $\{a\}$) and three inductive cases (union,
  concatenation, Kleene closure). A \textbf{regular expression} -- next lecture -- is the
  notation we use to name one.
\end{itemize}
\end{takeaways}

\section*{References for This Lecture}
\begin{thebibliography}{99}
\bibitem[DB]{db} Aho, A.~V., Lam, M.~S., Sethi, R., and Ullman, J.~D. \textit{Compilers:
Principles, Techniques, and Tools}. 2nd~ed. Pearson, 2006.
\bibitem[AP]{ap} Appel, A.~W. \textit{Modern Compiler Implementation in C/Java/ML}. Cambridge
University Press, 2004.
\end{thebibliography}

\end{document}
